\clearpage
\pdfbookmark[1]{Wirkung eines leuchtenden Punkts}{luminous-point}
\begin{center}\textbf{Wirkung eines leuchtenden Punkts}\end{center}
\noindent Diese mit KI erstellte Fassung in modernem Deutsch beruht auf dem deutschen Original \textit{Allgemeine Formeln für die Wirkung eines leuchtenden Punkts \(P\) auf einen Punkt \(p\)} in Carl Friedrich Gauss, \textit{Werke}, Band~V, S.~635--636, herausgegeben von der Königlichen Gesellschaft der Wissenschaften zu Göttingen. Siehe das \href{https://gdz.sub.uni-goettingen.de/id/PPN236006339}{Göttinger Digitalisat} und den zweiten Abdruck (1877), \href{https://archive.org/download/Werkecarlf05gausrich/page/n646}{S.~635} und \href{https://archive.org/download/Werkecarlf05gausrich/page/n647}{S.~636}.
\par\smallskip\noindent Zur optischen Terminologie siehe J. A. Fleming, \href{https://www.gutenberg.org/ebooks/71757}{\textit{Waves and Ripples in Water, Air, and Æther}}, zweite, überarbeitete Ausgabe, London 1912, Kap.~III, S.~105--106, sowie Marcel Indlekofer, Thomas Lauermann, Vincent Peikert und Raphael Straub, \href{https://theorie.physik.uni-konstanz.de/lsfuchs/lectures/ik30405/ik3optik1802.pdf#page=97}{\textit{Integrierter Kurs III -- Optik}}, Universität Konstanz, §1.10.2, S.~97. Das Skript beruht auf Vorlesungen von Georg Maret und Matthias Fuchs im Wintersemester 2004/2005.
\clearpage
